\section{Scaling Laws 与 LLM 的未来}

\subsection{Scaling Laws：收益递减但远未终结}

针对近期``Scaling Laws 已死''的论调，Hassabis 给出了明确的反驳：\textbf{Scaling Laws 的收益在递减，但远未触及天花板}（``diminishing returns but not dead''）。

\begin{importantbox}{Scaling Laws 的真实现状}
\begin{itemize}
    \item Scaling Laws \textbf{没有被突破}（have NOT been hit）
    \item 收益在递减，但曲线仍在下降
    \item 竞争的本质正在发生转变：从``比谁钱多''转向``比谁能想出新东西''
\end{itemize}
\textit{``收益递减但没死。真正的转变是竞争从'比谁钱多'变成'比谁能想出新东西'。''（Scaling: diminishing returns but not dead, competition shifting from ``who has more money'' to ``who can think of new things.''）}
\end{importantbox}

这个判断意味着：单纯堆算力的时代正在过去，但这并不意味着 Scaling 本身失效——而是进入了一个需要更多创造性思维的新阶段。算法创新、新的训练范式、更高效的架构设计将成为新的竞争焦点。

\subsection{算力：第一大瓶颈}

Hassabis 明确指出，\textbf{算力（Compute）是当前通往 AGI 的第一大瓶颈}。尽管竞争正在转向算法创新，但计算资源的规模仍然是前沿 AI 实验室之间拉开差距的关键因素。

\begin{warningbox}{LLM 不会被 Commoditize}
一个常见的误解是``LLM 最终会变成日用品，谁都能做''。Hassabis 对此表示明确反对：
\begin{itemize}
    \item LLM 不会被商品化（will NOT be commoditized）
    \item 前沿 AI 公司之间的差距\textbf{不是在缩小，而是在拉大}
    \item 只有前沿的 4 家左右的公司能持续推进 frontier 模型
\end{itemize}
原因在于：前沿模型需要的算力、数据、人才和算法 know-how 的门槛都在同时上升。
\end{warningbox}

\subsection{LLM 是 AGI 的基础组件}

尽管 LLM 有诸多短板，Hassabis 明确表示 LLM 将是未来 AGI 的\textbf{基础组件}（foundational component），而非被替代的过渡技术。关键不在于抛弃 LLM，而在于在 LLM 之上叠加缺失的能力：持续学习、记忆、规划和一致性推理。

\begin{knowledgebox}{从 LLM 到 AGI 的技术路径}
可以将 LLM 理解为 AGI 的``语言中枢''——它提供了强大的语言理解和生成能力，但 AGI 还需要：
\begin{itemize}
    \item 持续学习系统（类似海马体的记忆巩固）
    \item 层级规划模块（类似前额叶皮层的执行功能）
    \item 世界模型（对物理和社会世界的内部模拟）
    \item 一致性推理引擎（跨领域的逻辑一致性保证）
\end{itemize}
LLM 不会被替换，但会被大幅增强和扩展。
\end{knowledgebox}

\subsection{感知差距：短期高估、长期低估}

Hassabis 指出了一个重要的认知偏差：

\begin{importantbox}{AI 的感知差距}
\textit{``AI 在短期内被过度炒作，在长期（10 年尺度）被严重低估。''（AI overhyped short-term, underappreciated long-term.）}

这种感知差距导致了两个问题：
\begin{itemize}
    \item \textbf{短期}：过度的期望导致失望和``AI 寒冬''的恐慌
    \item \textbf{长期}：低估变革速度导致社会准备不足
\end{itemize}
正确的态度是：对近期进展保持务实，对长期影响保持敬畏。
\end{importantbox}

\subsection{本章小结}

Scaling Laws 远未终结，但竞争焦点正从资本密集型转向创新密集型。LLM 不会被商品化，前沿实验室的差距在拉大而非缩小。LLM 将作为 AGI 的基础组件被持续增强。公众对 AI 的感知存在``短期高估、长期低估''的系统性偏差。

%% ==================== Section 4 ====================
